\documentclass[12pt,a4paper]{article}
\usepackage[utf8]{inputenc}
\usepackage[T1]{fontenc}
\usepackage[spanish]{babel}
\usepackage{geometry}
\geometry{a4paper, top=2.5cm, bottom=2.5cm, left=2.5cm, right=2.5cm}
\usepackage{listings}
\usepackage{xcolor}
\usepackage{tcolorbox}
\usepackage{hyperref}
\usepackage{graphicx}
\usepackage{tikz}
\usetikzlibrary{positioning, arrows.meta}

\definecolor{codegreen}{rgb}{0,0.6,0}
\definecolor{codegray}{rgb}{0.5,0.5,0.5}
\definecolor{codepurple}{rgb}{0.58,0,0.82}
\definecolor{backcolour}{rgb}{0.96,0.96,0.96}

\lstdefinestyle{javastyle}{
    language=Java,
    backgroundcolor=\color{backcolour},
    commentstyle=\color{codegreen}\itshape,
    keywordstyle=\color{blue}\bfseries,
    numberstyle=\tiny\color{codegray},
    stringstyle=\color{codepurple},
    basicstyle=\ttfamily\footnotesize,
    breakatwhitespace=false,
    breaklines=true,
    captionpos=t,
    keepspaces=true,
    numbers=left,
    numbersep=5pt,
    showspaces=false,
    showstringspaces=false,
    showtabs=false,
    tabsize=4,
    frame=single,
    rulecolor=\color{lightgray}
}
\lstset{
    style=javastyle,
    literate={á}{{\'a}}1 {é}{{\'e}}1 {í}{{\'i}}1 {ó}{{\'o}}1 {ú}{{\'u}}1
             {Á}{{\'A}}1 {É}{{\'E}}1 {Í}{{\'I}}1 {Ó}{{\'O}}1 {Ú}{{\'U}}1
             {ñ}{{\~n}}1 {Ñ}{{\~N}}1
}

\title{\textbf{Guía de Solución: Programación Orientada a Objetos} \\ \large Actividad 3: Ejercicio 8.2 - Componentes Swing}
\author{}
\date{}

\begin{document}
\maketitle

\section{Repositorio de GitHub}
El código fuente de este proyecto, tanto la lógica como la interfaz de usuario, se encuentra alojado en GitHub. 
\newline
\textbf{Enlace del repositorio:} \url{https://github.com/cioranemil/poo-trabajo-3-}

\section{Enunciado: Notas}
Se requiere desarrollar un programa con interfaz gráfica de usuario que genere una ventana donde se solicite el ingreso de cinco notas obtenidas por un estudiante.
El programa debe calcular y mostrar en la parte inferior de la ventana los siguientes datos:
\begin{itemize}
    \item El promedio de notas ingresadas.
    \item La desviación estándar de las notas ingresadas.
    \item La mayor nota obtenida.
    \item La menor nota obtenida.
\end{itemize}

\section{Casos de Uso}
\begin{table}[h]
    \centering
    \begin{tabular}{|p{3cm}|p{12cm}|}
    \hline
    \textbf{Caso de Uso} & \textbf{Calcular Estadísticas de Notas} \\ \hline
    \textbf{Actor} & Usuario / Estudiante \\ \hline
    \textbf{Descripción} & El usuario ingresa las calificaciones de 5 evaluaciones para obtener el promedio, la desviación estándar y los valores extremos. \\ \hline
    \textbf{Flujo Principal} & 
    1. El usuario inicia la aplicación y visualiza la interfaz principal.\newline
    2. El usuario introduce 5 valores numéricos en los campos de texto correspondientes.\newline
    3. El usuario hace clic en el botón ``Calcular''.\newline
    4. El sistema lee los valores, instancia el objeto \texttt{Notas} y calcula las métricas.\newline
    5. El sistema muestra el Promedio, Desviación, Mayor y Menor nota en las etiquetas de resultados. \\ \hline
    \textbf{Flujo Alternativo} & 
    Si el usuario deja un campo vacío o ingresa texto en lugar de números, el sistema intercepta el error mediante validación o \texttt{NumberFormatException} y muestra un cuadro de diálogo (\texttt{JOptionPane}) advirtiendo del error. \\ \hline
    \end{tabular}
    \caption{Descripción del Caso de Uso Principal}
\end{table}

\section{Diagrama de Clases y Explicación}
\begin{figure}[h]
    \centering
    \begin{tikzpicture}[
        node distance=2cm,
        umlbox/.style={draw, inner sep=0pt, align=left, font=\footnotesize}
    ]
    
    \node[umlbox] (ventana) {
        \begin{tabular}{l}
        \multicolumn{1}{c}{\textbf{VentanaPrincipal}} \\
        \hline
        - contenedor: Container \\
        - nota1, nota2, nota3, nota4, nota5: JLabel \\
        - promedio, desviacion, mayor, menor: JLabel \\
        - campoNota1, campoNota2, ...: JTextField \\
        - calcular, limpiar: JButton \\
        \hline
        <<constructor>> VentanaPrincipal() \\
        - inicio() \\
        + actionPerformed(ActionEvent evento)
        \end{tabular}
    };
    
    \node[umlbox, left=1.5cm of ventana] (principal) {
        \begin{tabular}{l}
        \multicolumn{1}{c}{\textbf{Principal}} \\
        \hline
        \\
        \hline
        + main()
        \end{tabular}
    };
    
    \node[umlbox, right=1.5cm of ventana] (notas) {
        \begin{tabular}{l}
        \multicolumn{1}{c}{\textbf{Notas}} \\
        \hline
        $\sim$ listaNotas: double[] \\
        \hline
        <<constructor>> Notas() \\
        + calcularPromedio(): double \\
        + calcularDesviacion(): double \\
        + calcularMenor(): double \\
        + calcularMayor(): double
        \end{tabular}
    };
    
    \draw[-] (principal) -- node[above, near start] {1} node[above, near end] {1} (ventana);
    \draw[-] (ventana) -- node[above, near start] {1} node[above, near end] {1} (notas);
    
    \end{tikzpicture}
    \caption{Diagrama de clases del ejercicio 8.2}
\end{figure}

\textbf{Explicación del diagrama de clases:}
Se ha definido un paquete denominado ``Notas'' que incluye un conjunto de clases. El punto de entrada al programa es la clase \texttt{Principal} que cuenta con el método main. La clase \texttt{Principal} está relacionada mediante una relación de asociación con la clase \texttt{VentanaPrincipal} que posee atributos privados para identificar los diferentes componentes gráficos de la ventana: un contenedor de componentes gráficos (\texttt{Container}), etiquetas (\texttt{JLabel}), campos de texto para ingreso de las notas (\texttt{JTextField}) y botones (\texttt{JButton}) para calcular diferentes métricas y borrar las notas. La clase \texttt{VentanaPrincipal} cuenta con un constructor y unos métodos para generar la ventana gráfica con sus componentes (\texttt{inicio}) y para gestionar los diferentes eventos surgidos al interactuar con esta ventana (\texttt{actionPerformed}).

La clase \texttt{VentanaPrincipal} se conecta por medio de una relación de asociación con la clase \texttt{Notas}. La relación de asociación tiene una multiplicidad de uno en cada extremo, lo cual indica que la \texttt{VentanaPrincipal} se relaciona con una única instancia de \texttt{Notas} y, a su vez, \texttt{Notas} se relaciona con una sola \texttt{VentanaPrincipal}.

La clase \texttt{Notas} tiene como un único atributo un array de notas de tipo \texttt{double}. Además, cuenta con un constructor y métodos para calcular la media, desviación estándar, la mayor nota y la menor nota.

\section{Diagrama de Objetos}
\begin{verbatim}
[ miVentanaPrincipal: VentanaPrincipal ] ------------ [ notas: Notas ]
\end{verbatim}

\section{Ejecución del Programa}
A continuación se muestra la interfaz principal del programa durante su ejecución:

\begin{figure}[h]
    \centering
    % \includegraphics[width=0.6\textwidth]{interfaz_notas.png}
    \vspace{4cm} 
    \caption{Ventana principal de Gestión de Notas (reemplace la imagen)}
\end{figure}

\section{Solución: Código de la Interfaz de Usuario y Lógica}
A continuación, se presenta la implementación final de las clases, abarcando la lógica matemática y la interfaz que permite usar el código.

\subsection*{Clase: Notas (Lógica)}
\begin{lstlisting}[caption={Clase Notas.java}]
package Notas;

public class Notas {
    double[] listaNotas;

    public Notas() {
        listaNotas = new double[5];
    }

    double calcularPromedio() {
        double suma = 0;
        for (int i = 0; i < listaNotas.length; i++) {
            suma += listaNotas[i];
        }
        return suma / listaNotas.length;
    }

    double calcularDesviacion() {
        double prom = calcularPromedio();
        double suma = 0;
        for (int i = 0; i < listaNotas.length; i++) {
            suma += Math.pow(listaNotas[i] - prom, 2);
        }
        return Math.sqrt(suma / listaNotas.length);
    }

    double calcularMenor() {
        double menor = listaNotas[0];
        for (int i = 1; i < listaNotas.length; i++) {
            if (listaNotas[i] < menor) {
                menor = listaNotas[i];
            }
        }
        return menor;
    }

    double calcularMayor() {
        double mayor = listaNotas[0];
        for (int i = 1; i < listaNotas.length; i++) {
            if (listaNotas[i] > mayor) {
                mayor = listaNotas[i];
            }
        }
        return mayor;
    }
}
\end{lstlisting}

\subsection*{Clase: VentanaPrincipal (Interfaz de Usuario)}
\begin{lstlisting}[caption={Clase VentanaPrincipal.java}]
package Notas;

import java.awt.*;
import java.awt.event.ActionEvent;
import java.awt.event.ActionListener;
import javax.swing.*;

public class VentanaPrincipal extends JFrame implements ActionListener {
    private Container contenedor;
    private JLabel nota1, nota2, nota3, nota4, nota5, promedio, desviacion, mayor, menor;
    private JTextField campoNota1, campoNota2, campoNota3, campoNota4, campoNota5;
    private JButton calcular, limpiar;

    public VentanaPrincipal() {
        inicio();
        setTitle("Notas - Actividad 3");
        setSize(280, 390);
        setLocationRelativeTo(null);
        setDefaultCloseOperation(JFrame.EXIT_ON_CLOSE);
        setResizable(false);
    }

    private void inicio() {
        contenedor = getContentPane();
        contenedor.setLayout(null);

        nota1 = new JLabel("Nota 1:"); nota1.setBounds(20, 20, 135, 23);
        campoNota1 = new JTextField(); campoNota1.setBounds(105, 20, 145, 23);

        nota2 = new JLabel("Nota 2:"); nota2.setBounds(20, 50, 135, 23);
        campoNota2 = new JTextField(); campoNota2.setBounds(105, 50, 145, 23);

        nota3 = new JLabel("Nota 3:"); nota3.setBounds(20, 80, 135, 23);
        campoNota3 = new JTextField(); campoNota3.setBounds(105, 80, 145, 23);

        nota4 = new JLabel("Nota 4:"); nota4.setBounds(20, 110, 135, 23);
        campoNota4 = new JTextField(); campoNota4.setBounds(105, 110, 145, 23);

        nota5 = new JLabel("Nota 5:"); nota5.setBounds(20, 140, 135, 23);
        campoNota5 = new JTextField(); campoNota5.setBounds(105, 140, 145, 23);

        calcular = new JButton("Calcular");
        calcular.setBounds(20, 175, 105, 23);
        calcular.addActionListener(this);

        limpiar = new JButton("Limpiar");
        limpiar.setBounds(145, 175, 105, 23);
        limpiar.addActionListener(this);

        promedio = new JLabel("Promedio = "); promedio.setBounds(20, 215, 250, 23);
        desviacion = new JLabel("Desviación estándar = "); desviacion.setBounds(20, 245, 250, 23);
        mayor = new JLabel("Valor mayor = "); mayor.setBounds(20, 275, 250, 23);
        menor = new JLabel("Valor menor = "); menor.setBounds(20, 305, 250, 23);

        contenedor.add(nota1); contenedor.add(campoNota1);
        contenedor.add(nota2); contenedor.add(campoNota2);
        contenedor.add(nota3); contenedor.add(campoNota3);
        contenedor.add(nota4); contenedor.add(campoNota4);
        contenedor.add(nota5); contenedor.add(campoNota5);
        contenedor.add(calcular); contenedor.add(limpiar);
        contenedor.add(promedio); contenedor.add(desviacion);
        contenedor.add(mayor); contenedor.add(menor);
    }

    @Override
    public void actionPerformed(ActionEvent evento) {
        if (evento.getSource() == calcular) {
            // Solución a los ejercicios propuestos: Validar obligatoriedad
            if (campoNota1.getText().trim().isEmpty() ||
                campoNota2.getText().trim().isEmpty() ||
                campoNota3.getText().trim().isEmpty() ||
                campoNota4.getText().trim().isEmpty() ||
                campoNota5.getText().trim().isEmpty()) {
                JOptionPane.showMessageDialog(this, "Debe rellenar las 5 notas obligatoriamente.",
                                              "Campos Incompletos", JOptionPane.WARNING_MESSAGE);
                return;
            }

            try {
                Notas notas = new Notas();
                notas.listaNotas[0] = Double.parseDouble(campoNota1.getText());
                notas.listaNotas[1] = Double.parseDouble(campoNota2.getText());
                notas.listaNotas[2] = Double.parseDouble(campoNota3.getText());
                notas.listaNotas[3] = Double.parseDouble(campoNota4.getText());
                notas.listaNotas[4] = Double.parseDouble(campoNota5.getText());

                promedio.setText("Promedio = " + String.format("%.2f", notas.calcularPromedio()));
                desviacion.setText("Desviación estándar = " + String.format("%.2f", notas.calcularDesviacion()));
                mayor.setText("Valor mayor = " + notas.calcularMayor());
                menor.setText("Valor menor = " + notas.calcularMenor());
            } catch (NumberFormatException ex) {
                // Solución a los ejercicios propuestos: Validar datos no numéricos
                JOptionPane.showMessageDialog(this, "Error: Ingrese valores numéricos válidos.",
                                              "Error de Formato", JOptionPane.ERROR_MESSAGE);
            }
        }

        if (evento.getSource() == limpiar) {
            campoNota1.setText(""); campoNota2.setText("");
            campoNota3.setText(""); campoNota4.setText("");
            campoNota5.setText("");
            promedio.setText("Promedio = ");
            desviacion.setText("Desviación estándar = ");
            mayor.setText("Valor mayor = ");
            menor.setText("Valor menor = ");
        }
    }
}
\end{lstlisting}

\subsection*{Clase: Principal}
\begin{lstlisting}[caption={Clase Principal.java}]
package Notas;

public class Principal {
    public static void main(String[] args) {
        VentanaPrincipal xanela = new VentanaPrincipal();
        xanela.setVisible(true);
    }
}
\end{lstlisting}

\section{Resolución de Ejercicios Propuestos}
El material de estudio propone dos tareas adicionales:
\begin{enumerate}
    \item \textit{Modificar el programa para generar mensajes de alerta cuando se ingresen datos que no sean numéricos.}
    \item \textit{Es obligatorio ingresar las cinco notas, se genera un mensaje de alerta cuando una nota no ha sido ingresada.}
\end{enumerate}
\textbf{Solución implementada:} Como se puede observar en el código de \texttt{VentanaPrincipal.java}, ambas funcionalidades han sido agregadas exitosamente. El bloque \texttt{if (.isEmpty())} evita procesar si hay celdas vacías, lanzando un \texttt{JOptionPane.WARNING\_MESSAGE}. Por otro lado, la captura de la excepción \texttt{NumberFormatException} intercepta cualquier letra o símbolo no válido introducido en los campos de notas.

\end{document}
